\PassOptionsToPackage{sols}{handout}
\documentclass[main]{subfiles}
\setcounterrefbeforedot{chapter}{m-ws4-1}
\setcounterrefafterdot{section}{m-ws4-1}
\addtocounter{section}{-1}
\usepackage{solutions}

\begin{document}

\ifSubfilesClassLoaded{%
  \chaptermark{\chapterfour}
}{}

\section{Indeterminate Forms and L'Hôpital's Rule} \label{ws4-1}

\begin{problemonly}
    \begin{framed}

\textbf{Key Ideas}:

\begin{itemize}[topsep=0pt,partopsep=0pt,parsep=8pt,itemsep=0pt]

    \item You should understand what is meant by an ``indeterminate form.''
      (You should especially understand the difference between the terms
      ``indeterminate'' and ``undefined''.)
    
    \item You should be able to recognize the indeterminate forms
      ``$\frac{0}{0}$'' and ``$\frac{\infty}{\infty}$'' and understand
      why these are indeterminate. For today, we will limit ourselves
      to these forms.
    
    \item You should be able to evaluate limits involving the indeterminate
      forms ``$\frac{0}{0}$'' and ``$\frac{\infty}{\infty}$'' using L'Hôpital's
      Rule.


    \item You should recognize the form ``$0\cdot\infty$'' as an indeterminate
  form and have a strategy for tackling it. 

\end{itemize}
\end{framed}
\end{problemonly}

\begin{enumerate}
  \item
    Evaluate the following limits, using what you learned in Math Ma.  

    \note{In a limit like this first one, students usually can see that
      the quantity is ``getting big,'' but they forget to think about
      whether that means the limit is $\infty$ or $-\infty$.
      
      Matt note - we might see students trying to reason through these appealing to relative growth rates, so it might be good to remind them of when that actually applies.}
    \begin{enumerate}
 

    \item 
      \begin{problem}[\vfill]
        $\displaystyle\lim_{x \to 2} \frac{x^2-4}{x-2}$
      \end{problem}

      \begin{solution}
        As $x \to 2$, both the numerator and denominator approach 0, so we
        try to factor and cancel:
        \[ \lim_{x \to 2} \frac{x^2 - 4}{x - 2} = \lim_{x \to 2}
        \frac{\cancel{(x - 2)} (x + 2)}{\cancel{x - 2}} = \lim_{x \to 2} x
        + 2 = \boxed{4}.\]
      \end{solution}

    \item
      \begin{problem}[\vfill]
        $\displaystyle\lim_{x \to 2^-} \frac{x^2+x}{x-2}$
      \end{problem}

      \begin{solution}
        As $x \to 2^-$, $x^2 + x \to 4$ and $x - 2 \to 0$.  Since the
        denominator approaches 0 and the numerator does not, we must look
        at the sign of the denominator.  In this case, whenever $x$ is
        just a bit less than $2$, $x - 2$ is negative.  (We can write this
        as, ``As $x \to 2^-$, $x - 2 \to 0^-$''; that is, $x - 2$
        approaches 0 from below.)  Therefore, $\displaystyle \lim_{x \to
          2^-} \frac{x^2 + x}{x - 2} = \boxed{-\infty}$.  (If you plugged
        a value of $x$ just less than 2 into $\frac{x^2 + x}{x - 2}$,
        you'd be dividing a number close to 4 by a negative number that's
        very close to 0; the result would be very negative.)
      \end{solution}

      \item 
      \begin{problem}[\vfill]
          $\displaystyle\lim_{x \to 0}\frac{\sin(x)}{x}$
      \end{problem}

      \begin{solution}
          $\displaystyle\lim_{x \to 0}\frac{\sin(x)}{x} = 1$ (from section 3.5)
      \end{solution}

    \item
      \begin{problem}[\vfill]
        $\displaystyle\lim_{x \to \infty} \frac{5x^2+2x}{3x^2-x + 1}$
      \end{problem}

      \begin{solution}
        As $x \to \infty$, both the numerator and denominator $\to
        \infty$.  As you saw in Math Ma, we just need to think about the
        dominant (that is, fastest-growing) terms in the numerator and
        denominator; in the numerator, the fastest-growing term is $5x^2$;
        in the denominator, it's $3x^2$.  Therefore, \[ \lim_{x \to
          \infty} \frac{5x^2 + 2x}{3x^2 - x + 1} = \lim_{x \to \infty}
        \frac{5x^2}{3x^2} = \boxed{\frac{5}{3}}.\]
      \end{solution}


  \item
    \begin{problem}[\vfill\newpage]
      What is the limit $\displaystyle \lim_{h \to 0} \frac{f(3 + h) -
        f(3)}{h}$? 
    \end{problem}

    \begin{solution}
      This is the definition of the derivative $\boxed{f'(3)}$ (whose
      value will of course depend on what the function $f$ is).
    \end{solution}

  \end{enumerate}

\item  A few of the limits in {{\#1}} were ``of the form
  $\frac{0}{0}$,'' and one was ``of the form $\frac{\infty}{\infty}$.''
  ``$\frac{0}{0}$'' and ``$\frac{\infty}{\infty}$''
  are examples of what we call
  ``indeterminate forms.''

  \begin{enumerate}
  \item
    \begin{problem}
      What do we mean when we say ``$\frac{0}{0}$'' is an
      indeterminate form?
    \end{problem}

    \begin{solution}
      When we're looking at a limit $\displaystyle \lim_{x \to a}
      \frac{f(x)}{g(x)}$, just knowing that $\displaystyle \lim_{x \to a}
      f(x) = 0$ and $\displaystyle \lim_{x \to a} g(x) = 0$ does not tell
      us the value of $\displaystyle \lim_{x \to a} \frac{f(x)}{g(x)}$.  

      For example, in {{\#1}}, the limits $\displaystyle
      \lim_{x \to 2} \frac{x^2 - 4}{x - 2}$ and $\displaystyle \lim_{x \to
        0} \frac{\sin x}{x}$ were both $\frac{0}{0}$ limits, but they have
      different values.  After figuring out that they were $\frac{0}{0}$
      limits, we had to do more work to actually evaluate them.
    \end{solution}

    \pspace{\stretch{1}}

  \item
    \begin{problem}
      Come up with a couple of examples that show that
      ``$\frac{\infty}{\infty}$'' is also an indeterminate form.
    \end{problem}

    \begin{solution}
      Here are a few simple examples of $\frac{\infty}{\infty}$ limits:
      \begin{itemize}
      \item $\displaystyle \lim_{x \to \infty} \frac{x^2}{x} = \infty$.
      \item $\displaystyle \lim_{x \to \infty} \frac{x}{x} = 1$.
      \item $\displaystyle \lim_{x \to \infty} \frac{x}{x^2} = 0$.
      \end{itemize}
      The fact that they have different values shows that
      ``$\frac{\infty}{\infty}$'' is an indeterminate form.
    \end{solution}
  \end{enumerate}
  \pspace{\stretch{1}}

\item 
  \note{The goal of this problem is to motivate L'H\^opital's Rule.  (It's
    also fine to just tell them the rule.)}

  Let's look at the limit $\displaystyle \lim_{x \to 0}
  \frac{\sin x - 3x}{e^x -1}$.

  \begin{enumerate}

  \item

    \note{Here, I want them to find linear approximations for both
      $\sin x$ and $e^x$.  You might like to
      use a computer to actually zoom in on the two functions;
      \url{https://www.desmos.com/calculator} is good for this.}

    \begin{problem}[\stretch{3}]
      If you were to graph $\sin x - 3x$ and $e^x - 1$
      and zoom in around $x = 0$, what would the graphs resemble?  Can you
      approximate the graphs by something simpler near $x = 0$?
      Use tangent line approximations for $\sin x$ and $e^x$ at $x=0$.
    \end{problem}

    \begin{solution}

      \begin{itemize}
      \item  We've seen previously that $\sin x \approx x$ is the linear
      approximation of $\sin x$ when $x$ is around $0$. 
      (The derivative of $\sin x $ at $x=0$ is $1$, and $\sin 0 = 0$.)
      Thus, the graph of $\sin x - 3x$ will resemble the graph of 
      $x-3x=-2x$ around $x=0$.

      \item The derivative of $e^x$ at $x=0$ is $1$, and $e^0 = 1$,
      so $e^x \approx x+1$ around $x=0$. This means that the graph of
      $e^x - 1$ will resemble the graph of $x+1-1=x$ around $x=0$.
      \end{itemize}
    \end{solution}

  \item

    \note{Once we have $\sin x \approx x$ for $x$ near $0$ and
      $e^x \approx 1+x$ for $x$ near $0$, they should believe that
      \[
      \lim_{x \to 0} \frac{\sin x - 3x}{e^x - 1} = 
      \lim_{x \to 0} \frac{x - 3x}{1+x-1} = 
      \lim_{x \to 0} \frac{-2x}{x} = 
      -2.
      \]
      This suggests that the right rule for $\frac{0}{0}$ limits might be
      $\displaystyle \lim_{x \to a} \frac{f(x)}{g(x)} =
      \frac{f'(a)}{g'(a)}$.  Then I'll let them know what L'Hôpital's
      Rule actually says.  (I'll also let them know that it works for
      $\frac{\infty}{\infty}$ limits, because it's important to know
      the relative growth rates..) }

    \begin{problem}[\stretch{1.5}]
      Using your approximations, what do you think $\displaystyle \lim_{x
        \to 0} \frac{\sin x - 3x}{e^x - 1}$ is? Check on Desmos.
    \end{problem}

    \begin{solution}
      Since $\sin x - 3x \approx x -3x = -2x$ and $e^x-1 \approx x+1 -1 = x$ around $x=0$, we have
      
\[\frac{\sin
        (x) - 3x}{{e^x} - 1} \approx \frac{-2x}{x} = -2  \text{ for $x$ near $0$}.\]
      
      Therefore, a good guess is that $\displaystyle \lim_{x \to 0}
      \frac{\sin (x) - 3x}{e^{x} -1} = -2.$
    \end{solution}

  \end{enumerate}

\pspace{\newpage}

    \begin{framed}
      \textbf{L'Hôpital's Rule.}  If $\displaystyle \lim_{x \to a}
      \frac{f(x)}{g(x)}$ is a ``$\displaystyle \frac{0}{0}$'' or
      ``$\displaystyle \frac{\infty}{\infty}$'' type of limit, then
      $\displaystyle \lim_{x \to a} \frac{f(x)}{g(x)} = \lim_{x \to a}
      \frac{f'(x)}{g'(x)}$, if the latter limit exists or is $\pm \infty$.
      (This rule also works for one-sided limits.)

      {\Large\ding{43}} Remember that L'Hôpital's Rule only applies to
      limits of the form $\frac{0}{0}$ and $\frac{\infty}{\infty}$!

      \emph{Note: L'Hôpital's Rule depends heavily on all the work we've
      done learning how to take derivatives. It would not be possible
      to use this method had we not done that previous work.}
    \end{framed}

  \begin{enumerate}[resume]
  \item
    \begin{problem}[\stretch{1.5}]
      Use L'Hôpital's Rule to evaluate $\displaystyle \lim_{x
        \to 0} \frac{\sin x - 3x}{e^x - 1}$.
    \end{problem}

    \begin{solution}
      Since $\displaystyle \lim_{x
        \to 0} \frac{\sin x - 3x}{e^x - 1}$ is a 
        ``$\frac{0}{0}$'' limit, we can
      apply L'Hôpital's Rule:
      \begin{align*}
        \displaystyle \lim_{x
        \to 0} \frac{\sin x - 3x}{e^x - 1}
        &\lheq \lim_{x \to 0} \frac{\cos x - 3}{e^x}\\
        \intertext{As $x \to 0$, the numerator 
        $\cos x - 3 \to \cos(0) - 3 = -2$ and the denominator 
        $e^x \to 0$, so this limit is equal to}
        & \nolheq \frac{-2}{1} \\
        & \nolheq \boxed{-2}
      \end{align*}
    \end{solution}
  \end{enumerate}

\item 
  Evaluate the following limits.  \emph{Warning:} L'Hôpital's Rule does
  not apply to all of them!
  \begin{enumerate}
  \item
    \begin{problem}[100pt]
      $\displaystyle \lim_{x \to 1} \frac{x - 1}{\ln x}$      
    \end{problem}

    \begin{solution}
      As $x \to 1$, $x - 1 \to 0$ and $\ln x \to \ln 1 =
      0$, so we can use L'Hôpital's Rule:
      \begin{align*}
        \lim_{x \to 1} \frac{x - 1}{\ln x} 
        & \lheq \lim_{x \to 1} \frac{1}{1/x} \\
        \intertext{As $x \to 1$, $\frac{1}{x} \to 1$:}
        & \nolheq \boxed{1}
      \end{align*}
    \end{solution}

  \item

    \note{Need to use L'Hôpital's Rule twice here (or recognize the
      definition of a derivative after the first use).}

    \begin{problem}[\vfill]
      $\displaystyle \lim_{x \to 0} \frac{e^x - 1 - x}{x^2}$
    \end{problem}

    \begin{solution}
      This is a $\frac{0}{0}$ limit, so
      \begin{align*}
        \lim_{x \to 0} \frac{e^x - 1 - x}{x^2} & \lheq \lim_{x \to 0}
        \frac{e^x - 1}{2x} \\
        \intertext{This is again a $\frac{0}{0}$ limit:}
        & \lheq \lim_{x \to 0} \frac{e^x}{2} \\
        &= \boxed{\frac{1}{2}}
      \end{align*}
    \end{solution}

  \item

    \note{This is to make the point that L'Hôpital's Rule should not be
      used on limits that are not ``$\frac{0}{0}$'' or
      ``$\frac{\infty}{\infty}$''.  (If they try L'Hôpital's Rule,
      they'll get the wrong answer.)}

    \begin{problem}[\vfill\newpage]
      $\displaystyle \lim_{x \to 2} \frac{3x + 1}{x - 1}$
    \end{problem}

    \begin{solution}
      As $x \to 2$, $3x + 1 \to 7$ and $x - 1 \to 1$, so this limit is
      just equal to $\frac{7}{1} = \boxed{7}$.  (L'Hôpital's Rule does
      not apply, but we don't need it anyway!) 
    \end{solution}

  \item
    \note{If students work blindly here, they may end up using
      L'Hôpital's Rule over and over.  (After the first time, they get
      $\displaystyle \lim_{x \to \infty} \frac{1 / x}{x^{-2/3} / 3}$,
      which they might think of as a 0/0 limit.)  So, I'll use this to
      point out the value of simplifying before trying L'Hôpital's Rule
      again.}

    \begin{problem}[\vfill]
      $\displaystyle \lim_{x \to \infty} \frac{\ln x}{\sqrt[3]{x}}$
    \end{problem}

    \begin{solution}
      This is an $\frac{\infty}{\infty}$ limit:
      \begin{align*}
        \lim_{x \to \infty} \frac{\ln x}{\sqrt[3]{x}} & \lheq \lim_{x \to
          \infty} \frac{\frac{1}{x}}{\frac{1}{3} x^{-2/3}} \\
        & \nolheq \lim_{x \to \infty} \frac{3}{x^{1/3}} \\
        &= \boxed{0}
      \end{align*}
    \end{solution}
\end{enumerate}



\begin{enumerate}[start=5]

  \item

    \note{This is to remind students that $\frac{\infty}{0}$ is not
      indeterminate, but you do need to think about the sign of the
      denominator (and also to remind them about $\arcsin$).}

    \begin{problem}[\vfill]
      $\displaystyle \lim_{x \to 3^-} \frac{\ln (9 - x^2)}{\arcsin (x -
        3)}$
    \end{problem}

    \begin{solution}
      As $x \to 3^-$, $9 - x^2 \to 0^+$ (that is, when $x$ is just a bit
      less than $3$, $9 - x^2$ will be just a bit bigger than 0), so
      $\ln(9 - x^2) \to -\infty$.  Also, as $x \to 3^-$, $\arcsin(x - 3)
      \to 0^-$.  So, this limit is $\boxed{\infty}$. 
    \end{solution}

  \end{enumerate}


  \item  

  \note{Matt note - I would make sure to at least set up one of these limits for them before ending class, as there is one (fairly straightforward) problem on the PSET that is of this form. The point I try to repeatedly emphasize is that we need to use algebra to get a $"0 \cdot \infty"$ form into one that we can apply L'Hopital's to.}
  
  \begin{problem}[\nstretch{2}]
    Evaluate $\displaystyle \lim_{x \to 0^+} \sqrt{x} \ln x$.  What form
    is this limit?

    The strategy for dealing with a limit of this form is to rewrite the
    product $f(x) g(x)$ as a quotient; $f(x) g(x) = \frac{f(x)}{1 / g(x)}$
    or $f(x) g(x) = \frac{g(x)}{1 / f(x)}$.  (Both equations are true;
    often, one will be easier to work with than the other.)
  \end{problem}

  \begin{solution}
    $\displaystyle \lim_{x \to 0^+} \sqrt{x} = 0$ and $\displaystyle
    \lim_{x \to 0^+} \ln x = -\infty$.  Therefore, this is a ``$0 \cdot
    \infty$'' type of limit, so we first use algebra to rewrite the
    product as a fraction.  There are two ways we could do this:
    \begin{align*}
      \sqrt{x} \ln x &= \frac{\ln x}{1 / \sqrt{x}} & \text{or} &&
      \sqrt{x} \ln x &= \frac{\sqrt{x}}{1 / \ln x}
    \end{align*}
    These two equations are both correct, so we can use either of them.
    The first one looks computationally a little simpler because we know
    that we can use exponentiation rules to rewrite $1 / \sqrt{x}$ as
    $x^{-1/2}$.  So, let's use the left one and rewrite $\displaystyle
    \sqrt{x} \ln x = \frac{\ln x}{x^{-1/2}}$.  Now, $\displaystyle
    \lim_{x \to 0^+} \frac{\ln x}{x^{-1/2}}$ is a
    ``$\frac{\infty}{\infty}$'' type of limit, so we can use
    L'Hôpital's Rule to evaluate it:
    \begin{align*}
      \lim_{x \to 0^+} \frac{\ln x}{x^{-1/2}} &
      \lheq \lim_{x \to 0^+}
      \frac{\frac{1}{x}}{ - \frac{1}{2} x^{-3/2}} \\
      \intertext{When written this way, this limit still looks like a
        ``$\frac{\infty}{\infty}$'' type of limit; the key is to use
        algebra to rewrite it:} %
      & \nolheq \lim_{x \to 0^+} -2 \cdot \frac{1}{x} \cdot x^{3/2} \\
      & \nolheq \lim_{x \to 0^+} -2 \sqrt{x} \\
      \intertext{This is a limit we can evaluate simply by plugging in
        $x = 0$:}  %
      &= \boxed{0}
    \end{align*}      
  \end{solution}

\item 

  \note{Just more practice recognizing and working with $0 \cdot \infty$
    limits.} 

  Evaluate the following limits.  
  \begin{enumerate}
  \item
    \begin{problem}[\vfill]
      $\displaystyle \lim_{x \to-\infty} x e^{-2x}$
    \end{problem}

    \begin{solution}
      As $x \to -\infty$, $e^{-2x} \to \infty$, so this limit has form
      $-\infty \cdot \infty$, and it's just $\boxed{-\infty}$.  (As $x$
      becomes very negative, $x e^{-2x}$ involves multiplying a very
      negative number with a huge positive number, which gives an
      extremely negative result.)
    \end{solution}

  \item
    \begin{problem}[\vfill]
      $\displaystyle \lim_{x \to -\infty} x^2 e^{x}$
    \end{problem}

    \begin{solution}
      As $x \to -\infty$, $x^2 \to \infty$ and $e^x \to 0$.  So, this is
      a $0 \cdot \infty$ limit, and we need to rewrite it and use
      L'H\^opital's Rule:
      \begin{align*}
        \lim_{x \to -\infty} x^2 e^x &\nolheq \lim_{x \to -\infty}
        \frac{x^2}{e^{-x}} \\
        \shortintertext{This is now an $\frac{\infty}{\infty}$ limit:}
        & \lheq \lim_{x \to -\infty} \frac{2x}{-e^{-x}} \\
        \shortintertext{This is another $\frac{\infty}{\infty}$ limit:}
        & \lheq \lim_{x \to - \infty} \frac{2}{e^{-x}} \\
        \intertext{Now, the numerator is 2 while the denominator $\to
          \infty$, so the limit is equal to}
        & \nolheq \boxed{0}
      \end{align*}
    \end{solution}

  \end{enumerate}

\end{enumerate}

  \begin{framed}
    \textbf{$0 \cdot \infty$ limits.}  The strategy for dealing
    with a limit $\displaystyle \lim_{x \to a} f(x) g(x)$ of the
    indeterminate form $0 \cdot \infty$ is to rewrite the product $f(x)
    g(x)$ as a quotient; $f(x) g(x) = \frac{f(x)}{1 / g(x)}$ or $f(x)
    g(x) = \frac{g(x)}{1 / f(x)}$.  (Both equations are true; often, one
    will be easier to work with than the other.)
  \end{framed}

\end{document}